\documentclass[11pt]{article}

% Compile with: pdflatex E26_regularized_curvature_L2_bound.tex
\usepackage[T1]{fontenc}
\usepackage{lmodern}
\usepackage[margin=1in]{geometry}
\usepackage{amsmath,amssymb,amsthm,mathtools}
\usepackage{microtype}
\usepackage[colorlinks=true,linkcolor=blue,urlcolor=blue]{hyperref}

\newtheorem{theorem}{Theorem}
\newtheorem{corollary}{Corollary}
\theoremstyle{remark}
\newtheorem{remark}{Remark}
\newtheorem{counterexample}{Counterexample}
\DeclareMathOperator*{\essinf}{ess\,inf}

\title{E26: A Conditional $L^2$ Estimate for the\\
Regularized Curvature Operator}
\author{}
\date{}

\begin{document}
\maketitle

\begin{abstract}
For $\varepsilon>0$ and $\alpha\geq0$, consider the curvature
expression
\[
\Theta_\varepsilon[\psi]
=-\frac{\Delta\psi}
{\psi+\varepsilon e^{-\alpha|\psi|^2}}.
\]
We prove the $L^2$ bound asserted in lemma E26 under a precise,
uniform lower bound on the modulus of the denominator. We also establish
two explicit sufficient conditions and exhibit a smooth counterexample
to the unconditional version of the estimate. In particular,
$\varepsilon>0$ alone does not prevent denominator zeros for signed
real or complex fields.
\end{abstract}

\section{Definitions and main result}
Let $\Omega\subseteq\mathbb R^d$ be open, with $d\geq1$.
Fix $\varepsilon>0$ and $\alpha\geq0$ and take
$\psi\in H^2(\Omega;\mathbb C)$. Define
\begin{equation}\label{eq:D}
D_\varepsilon[\psi](x)
:=\psi(x)+\varepsilon e^{-\alpha|\psi(x)|^2}.
\end{equation}
Where $D_\varepsilon[\psi]\neq0$, let
\begin{equation}\label{eq:Theta}
\Theta_\varepsilon[\psi](x)
:=-\frac{\Delta\psi(x)}{D_\varepsilon[\psi](x)}.
\end{equation}
Pointwise assertions about Sobolev representatives are always
understood almost everywhere.

\begin{theorem}[E26: conditional $L^2$ estimate]\label{thm:E26}
Assume that there is a constant $\delta>0$ for which
\begin{equation}\label{eq:H}
\essinf_{x\in\Omega}\bigl|D_\varepsilon[\psi](x)\bigr|
\geq\delta.
\end{equation}
Then $\Theta_\varepsilon[\psi]\in L^2(\Omega;\mathbb C)$ and
\begin{equation}\label{eq:E26}
\boxed{
\bigl\|\Theta_\varepsilon[\psi]\bigr\|_{L^2(\Omega)}
\leq\frac{1}{\delta}\,
\bigl\|\Delta\psi\bigr\|_{L^2(\Omega)}.
}
\end{equation}
The estimate is valid for every $d\geq1$; no additional boundary
condition is needed for this inequality.
\end{theorem}

\begin{proof}
Since $\psi\in H^2(\Omega;\mathbb C)$, every weak second partial
derivative of $\psi$ belongs to $L^2(\Omega)$, and hence
\[
\Delta\psi=\sum_{j=1}^{d}\partial_j^2\psi
\in L^2(\Omega;\mathbb C).
\]
The expression $D_\varepsilon[\psi]$ is measurable. Hypothesis
\eqref{eq:H} implies that it is nonzero almost everywhere and that
its reciprocal is essentially bounded:
\begin{equation}\label{eq:reciprocal}
\left\|\frac{1}{D_\varepsilon[\psi]}\right\|_{L^\infty(\Omega)}
\leq\frac{1}{\delta}.
\end{equation}
The product of an $L^\infty$ function and an $L^2$ function belongs
to $L^2$. In particular, by H\"older's inequality,
\begin{align*}
\bigl\|\Theta_\varepsilon[\psi]\bigr\|_{L^2(\Omega)}
&=\left\|\frac{\Delta\psi}{D_\varepsilon[\psi]}
\right\|_{L^2(\Omega)}\\
&\leq\left\|\frac{1}{D_\varepsilon[\psi]}
\right\|_{L^\infty(\Omega)}
\bigl\|\Delta\psi\bigr\|_{L^2(\Omega)}\\
&\leq\frac{1}{\delta}\,
\bigl\|\Delta\psi\bigr\|_{L^2(\Omega)}.
\end{align*}
This proves \eqref{eq:E26}. Equivalently, one can calculate directly:
\begin{align*}
\bigl\|\Theta_\varepsilon[\psi]\bigr\|_{L^2(\Omega)}^2
&=\int_\Omega
\frac{|\Delta\psi(x)|^2}{|D_\varepsilon[\psi](x)|^2}\,dx\\
&\leq\delta^{-2}\int_\Omega|\Delta\psi(x)|^2\,dx\\
&=\delta^{-2}\bigl\|\Delta\psi\bigr\|_{L^2(\Omega)}^2.
\end{align*}
Taking nonnegative square roots gives the same conclusion.
\end{proof}

\section{Sufficient denominator bounds}

\begin{corollary}[Fields bounded away from zero]\label{cor:modulus}
If $|\psi(x)|\geq m>\varepsilon$ almost everywhere on $\Omega$, then
\begin{equation}\label{eq:modulus}
\bigl\|\Theta_\varepsilon[\psi]\bigr\|_{L^2(\Omega)}
\leq\frac{1}{m-\varepsilon}\,
\bigl\|\Delta\psi\bigr\|_{L^2(\Omega)}.
\end{equation}
\end{corollary}

\begin{proof}
As $\alpha\geq0$, one has
$0<\varepsilon e^{-\alpha|\psi|^2}\leq\varepsilon$.
The reverse triangle inequality therefore yields
\[
|D_\varepsilon[\psi]|
\geq|\psi|-\varepsilon e^{-\alpha|\psi|^2}
\geq m-\varepsilon>0.
\]
Apply Theorem~\ref{thm:E26} with $\delta=m-\varepsilon$.
\end{proof}

\begin{remark}
On $\Omega=\mathbb R^d$, an $L^2$ function cannot satisfy
$|\psi|\geq m>0$ almost everywhere. Thus
Corollary~\ref{cor:modulus} is most naturally used on finite-measure
domains. The main theorem itself applies equally well on unbounded
domains whenever its hypothesis is satisfied.
\end{remark}

\begin{corollary}[Nonnegative real fields]\label{cor:nonnegative}
Suppose $\psi\in H^2(\Omega;\mathbb R)$ and $\psi\geq0$ almost
everywhere. For every fixed $\varepsilon>0$ and $\alpha\geq0$, set
\begin{equation}\label{eq:deltaeps}
\delta_\varepsilon
:=\min\left\{\frac{\varepsilon}{2},\,
\varepsilon e^{-\alpha\varepsilon^2/4}\right\}>0.
\end{equation}
Then
\begin{equation}\label{eq:nonnegative}
\bigl\|\Theta_\varepsilon[\psi]\bigr\|_{L^2(\Omega)}
\leq\frac{1}{\delta_\varepsilon}\,
\bigl\|\Delta\psi\bigr\|_{L^2(\Omega)}.
\end{equation}
\end{corollary}

\begin{proof}
For almost every $x$, write $r=\psi(x)\geq0$.
If $r\geq\varepsilon/2$, then
\[
r+\varepsilon e^{-\alpha r^2}\geq\varepsilon/2.
\]
If $0\leq r<\varepsilon/2$, the assumption $\alpha\geq0$ yields
\[
r+\varepsilon e^{-\alpha r^2}
\geq\varepsilon e^{-\alpha\varepsilon^2/4}.
\]
Thus $D_\varepsilon[\psi]\geq\delta_\varepsilon$ almost
everywhere. The assertion follows from Theorem~\ref{thm:E26}.
\end{proof}

\section{Counterexample without denominator separation}

\begin{counterexample}[Positive regularization is not enough]
Choose $\Omega=(-1,1)$, take any $\varepsilon>0$, set $\alpha=0$,
and define the smooth real-valued function
\[
\psi(x)=x^2-\varepsilon.
\]
Then $\psi\in H^2((-1,1))$ and
\[
D_\varepsilon[\psi](x)=x^2,
\qquad
\Delta\psi(x)=\psi''(x)=2.
\]
Although the denominator vanishes only at $x=0$, the quotient is not
square-integrable:
\[
\Theta_\varepsilon[\psi](x)=-\frac{2}{x^2}
\quad(x\ne0),
\qquad
\int_{-1}^{1}|\Theta_\varepsilon[\psi](x)|^2\,dx
=4\int_{-1}^{1}|x|^{-4}\,dx=+\infty.
\]
In contrast,
\[
\bigl\|\Delta\psi\bigr\|_{L^2((-1,1))}^2
=\int_{-1}^{1}4\,dx=8<\infty.
\]
Consequently, the E26 estimate is \emph{false} for unrestricted
signed real fields (and hence for unrestricted complex fields) if
one assumes only $\varepsilon>0$ and $\psi\in H^2$.
\end{counterexample}

\section{Scope and interpretation}
Hypothesis~\eqref{eq:H} is a bound on the \emph{modulus of the full
denominator}. It does not follow from the mere positivity of
$\varepsilon$. Nor may the constant $\delta$ be replaced by
$\varepsilon$ without additional restrictions on $\psi$.

Theorem~\ref{thm:E26} is a pointwise-in-$\psi$ multiplication estimate,
not an estimate for the difference of operators evaluated at two
different fields. It establishes neither local Lipschitz continuity
nor global existence or uniqueness for an associated evolution
problem. Furthermore, the theorem supplies no uniform-in-$\varepsilon$
estimate as $\varepsilon\downarrow0$ unless a positive denominator
bound independent of $\varepsilon$ is separately established.

\end{document}
